We site parks, water bodies and cool pavement on synthetic 50\,m city grids. Each instance is a
constrained binary program, encoded as a \QUBO{} or \HUBO{} and scored on exact saturating cooling
physics against a HiGHS-certified optimum. A slack-bit penalty \QUBO{} with a corrected weight has
the MILP optimum as its ground state on all 11 exhaustively enumerated instances. The second-order
Taylor \QUBO{} loses 10.3\% of the achievable cooling at 300\,m reach, where the cubic \HUBO{} loses
at most 0.02\% and a least-squares-fitted \QUBO{} at most 0.42\%. The final annealing gap scales as
$1/\Lam$ with the penalty weight. Constraint-preserving Dicke-XY QAOA reaches a mean $\Popt$ of
0.43 at depth 6, against 1.00 for feasible-space simulated annealing. Path-integral Monte Carlo
restricted to the feasible set passes a pre-committed comparability bar (benefit 0.981), but it
made 3.6$\times$ more proposals than feasible-space annealing; at equal proposals it scores 0.983
against 0.988. HiGHS certifies every hard-family instance in 0.77\,s on average. Encodings and
constraint handling, not the choice of annealer, decide the outcome. We observe no quantum or
quantum-inspired advantage.
